\section{Central solenoid and poloidal field coils}
\label{sec:cspf}
% NB: the CS/PF solvers are not yet in this repository
% (MADE-2.0/matlab contains only the TF chain). Import them before
% writing this section so equations match the code.

\subsection{Central solenoid}
% Input: flux swing (from scenario: ramp-up + flat-top + Ejima), CS
% radius/height from inboard build, number of modules, cycles.
% Physics: B_peak from flux and geometry; conductor via the shared
% module; jacket/structure sized on hoop + vertical (separating) loads,
% Tresca and fatigue (cyclic loading -> crack growth / allowable
% reduction). Pre-compression structure.
\todo{equations and design flow from the CS solver.}

\subsection{Poloidal field coils}
% Input: PF currents (per scenario snapshot) and positions from the
% free-boundary equilibrium. Field on coil from Biot-Savart of the full
% CS+PF+plasma set (or envelope). Conductor via shared module;
% rectangular WP; hoop stress in the jacket (ring); support loads.
\todo{equations and design flow from the PF solver.}

\subsection{Field on coil for axisymmetric sets}
% Common sub-module: field of thick axisymmetric coils (elliptic
% integrals / filament discretisation) -> peak field per coil and
% per grade. Also the loads (radial hoop, vertical) per coil.
% This is the sub-module re-used for mirrors (section 6).
\todo{describe the method used in the code (xbr.m/xbz.m?).}
